%%%%%%%%%%%%%%%%%%%%%%%%
% $Autor: MansourAhmadyPhoulady, RanjeetsingTawar, HemanthPrabhakaran$
% $Pfad: Car_Licence_Plate_Identification_with_a_JetsonNano$
% $Version: 4.0.0 $
%
% !TeX encoding = utf8
%
%%%%%%%%%%%%%%%%%%%%%%%%
\chapter{Introduction}
In present day scenario the security and authentication is very much needed to make a safety world.Beside all security one vital issue is recognition of number plate from the car for authorization. In the busy world everything cannot be monitor by a human, so automatic license plate recognition is one of the best application for authorization without involvement of human power.\\
The suggested solution belongs to the category of challenges categorized as "image processing," which can be tackled using Convolutional Neural Networks (CNNs), a subtype of Artificial Neural Networks. CNNs are chiefly designed for intricate image-based pattern recognition tasks, owing to their precise yet uncomplicated architecture. This solution comprises two distinct phases to determine whether a vehicle may be granted parking privileges. In the initial stage, the system undergoes training and testing using a substantial dataset to identify various segments of the license plate and accurately extract the pertinent information. Following this training, the system is deployed onto the Jetson Nano. In the subsequent stage, the Jetson Nano's onboard camera captures the license plate information from the car's image. If the recognition process is successful, a green indicator illuminates, signifying a valid license plate recognition. Conversely, if recognition errors occur or if the identified license plate is not found in the database, a red indicator is activated to signal that parking is not allowed for the vehicle.\cite{Sandha:2017}\\ 

\section{Challenges}
Recognizing vehicle plates has historically presented several challenges. These challenges include:\cite{Awalgaonkar:2021}

\begin{itemize}
	\item Variability in Plate Design: License plates come in a wide range of designs and formats, with variations in colors, fonts, and sizes. This variability can make it difficult to create a single, universal recognition system.
	\item Speed of Vehicles: In scenarios like toll booths or fast-moving traffic, the speed of vehicles can make it challenging to capture license plates quickly and accurately.
	\item Varying Plate Sizes: Different countries and regions have license plates of varying sizes, and this adds complexity to recognition systems that need to adapt to these differences.
	\item Foreign License Plates: In areas with international traffic, systems must be capable of recognizing a wide range of license plate designs from different countries.
	\item Processing Power: Efficiently processing images and running recognition algorithms in real-time, especially in scenarios with a large volume of vehicles, requires substantial computational power.
\end{itemize}

\section{Proposed Solution}
Recognizing vehicle plates presents various challenges, and YOLO (You Only Look Once) can assist in addressing some of these challenges\cite{Sung:2020}. These challenges encompass:

\begin{itemize}
	\item Variability in Plate Design: YOLO can aid in localizing and identifying license plates. However, the recognition of specific design elements, such as characters, colors, and fonts, may necessitate additional processing steps, such as Optical Character Recognition (OCR).
	
	\item Speed of Vehicles: YOLO is specifically designed for real-time object detection, making it well-suited for scenarios involving fast-moving vehicles. Nevertheless, the accuracy of license plate capture at high speeds is contingent on the camera setup and image quality.
	
	\item Varying Plate Sizes: YOLO exhibits adaptability to objects of varying sizes, including license plates, rendering it versatile for deployment in regions with different plate dimensions.
	
	\item Foreign License Plates: YOLO can be trained to recognize a diverse array of license plate designs, making it an apt choice for areas with international traffic.
	
	\item Processing Power: YOLO is renowned for its real-time processing capabilities. However, its actual performance is influenced by the hardware employed. High-performance GPUs are often necessary for achieving real-time processing.
\end{itemize}


\section{Report Structure}

\begin{enumerate}
	\item \textbf{Introduction}: Provides an overview of the project or subject matter.
	\item \textbf{Challenges}: Identifies and discusses the obstacles or difficulties faced in the project.
	\item \textbf{Data Version Control}: Discusses methods and best practices for managing and versioning data.
	\item \textbf{Data Description}: Describes the characteristics and attributes of the data used in the project.
	\item \textbf{Hardware Description}: Provides information about the hardware infrastructure utilized in the project.
	\item \textbf{Development Environment and Software Description}: Describes the software tools and environment used for development.
	\item \textbf{Data Mining}: Discusses techniques and processes for extracting insights and patterns from data.
	\item \textbf{KDD Process}: Explores the Knowledge Discovery in Databases (KDD) process for extracting useful information from data.
	\item \textbf{Development Folder Structures}: Describes the organization and structure of folders and directories used in development.
	\item \textbf{Document Developer}: Discusses the documentation requirements and practices for developers.
	\item \textbf{Programming Concepts}: Covers fundamental concepts and principles of programming relevant to the project.
	\item \textbf{Requirements.txt}: Specifies the dependencies and requirements for the project, typically in a text file format.
	\item \textbf{Software Tests}: Discusses the testing strategies and methodologies employed in the project.
	\item \textbf{Deployment}: Covers the process of deploying the software or solution into a production environment.
	\item \textbf{Monitoring}: Describes the methods and tools used for monitoring the performance and behavior of the deployed system.
	\item \textbf{Program Flowchart}: Provides a visual representation of the flow of the program or system logic.
	\item \textbf{Bill of Materials}: Lists the components and resources required for the project, typically used in manufacturing contexts.
	\item \textbf{Numpy Package}: Discusses the features and usage of the NumPy library for numerical computing in Python.
	\item \textbf{Pandas Package}: Discusses the features and usage of the Pandas library for data manipulation and analysis in Python.
	\item \textbf{Matplotlib Package}: Discusses the features and usage of the Matplotlib library for data visualization in Python.
	\item \textbf{Conclusion}: Summarizes the findings, outcomes, and conclusions drawn from the project.
	\item \textbf{Data Sheets}: Provides detailed documentation or specifications about the datasets used in the project.
\end{enumerate}